\documentclass[11pt]{amsart}
\ifdefined\pdftrailerid\pdftrailerid{}\fi
\usepackage[margin=1.2in]{geometry}
\usepackage{amsmath,amssymb,amsthm,amscd,mathtools}
\usepackage[colorlinks=true,linkcolor=blue,citecolor=blue,urlcolor=blue]{hyperref}
\hypersetup{pdftitle={The first variation and the one-mark coefficient of circle exchange for balanced local curves},pdfauthor={Bernd Johannes Wuebben}}
\setlength{\emergencystretch}{2em}
\newtheorem{theorem}{Theorem}[section]
\newtheorem{proposition}[theorem]{Proposition}
\newtheorem{lemma}[theorem]{Lemma}
\newtheorem{corollary}[theorem]{Corollary}
\theoremstyle{definition}\newtheorem{example}[theorem]{Example}
\newcommand{\Iorig}{\mathcal I_m^{\mathrm{orig}}}
\newcommand{\Irot}{\mathcal I_m^{\mathrm{rot}}}
\newcommand{\OO}{\mathcal O}
\newcommand{\AAo}{\mathbb A}
\newcommand{\CC}{\mathbb C}
\newcommand{\ZZ}{\mathbb Z}
\newcommand{\PP}{\mathbb P}
\DeclareMathOperator{\Pic}{Pic}
\DeclareMathOperator{\Hilb}{Hilb}
\DeclareMathOperator{\End}{End}
\DeclareMathOperator{\rk}{rk}
\title[Refined circle exchange for balanced local curves]{The first variation and the one-mark coefficient of circle exchange for balanced local curves}
\author{Bernd Johannes Wuebben}
\subjclass[2020]{14N35, 14C35, 14H60, 05E05}
\keywords{Stable pairs, local curves, equivariant $K$-theory, Picard varieties, Macdonald operators}
\date{November 26, 2026}
\begin{document}
\begin{abstract}
For a balanced local curve of positive genus we prove a first-order
circle-exchange identity for the two stable-pair indices over its
Picard variety. After their common determinant normalization, the
indices agree modulo $(1-\kappa)^2$ in integral family $K$-theory
in every curve degree and both counting gradings, where $\kappa$ is
the threefold volume character. At positive exchanged grading the
first variation is a sum over hooks of Adams operations on the
derived pushforwards of powers of the Poincar\'e line twisted by
powers of the canonical bundle. With $\kappa$ unrestricted, the
entire coefficient of exchanged grading one agrees separately for
each Young diagram and is given by a sum over its removable boxes.
The first comparison deforms the graded moment equations to walls
with symmetric powers of the curve as fixed loci. The second
evaluates a product of two Macdonald currents and descends the
resulting disk identity to the Picard family.
\end{abstract}
\maketitle
\section{The two family indices}
Circle exchange compares the stable-pair theories on the fixed threefolds of two different fiber rotations of a fivefold. Let $C$ be a smooth projective complex curve of genus $g=b+1\ge1$,
let $B=\Pic^b(C)$, and let $L$ be the Poincar\'e line on $C\times B$
trivialized at a chosen point of $C$. Put $K=K_C$, $M=K/L$ and
$D(F)=R\pi_*F$ for the projection $\pi:C\times B\to B$. Products
and quotients of lines denote tensor products and duals. For a
threefold whose normal bundle over $C$ is $L\oplus M$, the equivariant
stable-pair index can be compared with that of the product
$C\times\AAo^2$. The comparison varies $L$ over $B$ on both sides.

Consider a double curve in genus two. At the minimum Euler grading
the full normal torus fixes two pairs over each $L$, the ribbons in
the two normal directions. They have Euler characteristic $-3$. Write $q$ for the original Euler grading, $t$
for the exchanged grading, and $\delta=1-\kappa$ for the change in
the volume character. At $q^{-3}t^2$ the coefficient of $\delta$ is
$-6\psi^2 D(L)$, where $\psi^j$ denotes the integral Adams operation.
This class has rank zero. Its Chern character in rational singular
cohomology is $12\theta$, where $\theta$ is the principal polarization
of $\Pic^1(C)$. Indeed Grothendieck--Riemann--Roch in cohomology gives
$\operatorname{ch}D(L)=-\theta$.
Passing to the numerical index of a single target loses this coefficient.
The calculation at $q^{-3}$ raises two questions. Does the variation
agree at every Euler grading and in arbitrary curve degree? How much
of the comparison persists when $\kappa$ is free?

The two rotations act on the ordered summands of the fivefold
\[
 \operatorname{Tot}_{C\times B}
 (tL\oplus(\kappa/t)M\oplus q\OO\oplus(c/q)\OO),\qquad
 c=\kappa^{-1},\quad \sigma^2=\kappa .
\]
The original threefold uses the first two summands, the rotated one
the last two. A stable pair is a section of a pure one-dimensional
sheaf with proper support and zero-dimensional cokernel. Its degree
$m$ is the multiplicity over $C$. In the original index a pair of
Euler characteristic $j$ is counted with $(-q)^j$. In the rotated
index the corresponding monomial is $(-t)^k$, where
$k=\chi(\mathcal F)+mb$. We call $[t]$ the one-mark coefficient: on the product threefold a single box of the monomial quotient acquires a divisor of length one. The modified virtual structure sheaf uses
the ordered external lines in each presentation, as specified in
Section~\ref{sec:geometry}. Denote the resulting pushforwards to $B$
by $\Iorig,\Irot$. In a partition $\lambda$ write $a_x,\ell_x,h_x$
for the arm, leg and hook length $1+a_x+\ell_x$ of a box $x$. Set
\[
 Z_\lambda(q)=\prod_{x\in\lambda}[(1-q^{h_x})(1-q^{-h_x})]^b,\qquad
 D_m=bm(m+1)/2.
\]

\begin{theorem}\label{thm:first}
For every $m\ge1$, both $\sigma^{mb}\Iorig$ and $\sigma^{mb}\Irot$
belong to $q^{-D_m}K^0(B)[\kappa^{\pm1}][[q,t]]$ and are congruent
modulo $(1-\kappa)^2$. For $k>0$ the common first variation is
\begin{equation}\label{eq:first}
 [\delta t^k]\mathcal I_m
 =-\sum_{\lambda\vdash m}Z_\lambda(q)
 \sum_{\substack{x\in\lambda\\h_x\mid k}}
 \psi^{k/h_x}D(L^{h_x}K^{-\ell_x}).
\end{equation}
The formula holds after restriction to any orbit of diagrams under
transposition. Here the positive-$t$ coefficient is divided by
$1-\kappa$ before specializing $\sigma=1$. At $t^0$, on such an
orbit, put $z_j=[q^j]\sum_\lambda Z_\lambda(q)$. The coefficient of
$q^j$ in either unnormalized character is
$z_j(1+j(\sigma-1))$ modulo $(\sigma-1)^2$.
\end{theorem}

For the second comparison set $u=q$, $v=c/u$, $p=K$, $s=p/c$,
$R(z)=(1-sz)/(1-z)$ and $d_{ij}=u^iv^j$, with $(0,0)$ the
northwest box, $i$ increasing to the right and $j$ downwards. Denote the removable and addable boxes of $\lambda$
by $\mathcal R_\lambda$ and $\mathcal A_\lambda$. Define
\begin{align}
 Z_\lambda(q,\kappa)
 &=(-1)^{mb}q^{-b\sum_x h_x}\kappa^{-b\sum_x(\ell_x+1/2)}
 \prod_x(1-\kappa^{\ell_x}q^{h_x})^b
         (1-\kappa^{\ell_x+1}q^{h_x})^b,\label{eq:empty}\\
 G_\lambda
 &=R(u)R(v)\sum_{r\in\mathcal R_\lambda}
 \frac{\prod_{h\in\mathcal R_\lambda\setminus\{r\}}R(d_r/d_h)}
 {\prod_{a\in\mathcal A_\lambda}R(cd_r/d_a)}.\label{eq:corners}
\end{align}

\begin{theorem}\label{thm:mark}
At unrestricted volume, for every $m\ge1$, the contributions of each
$\lambda\vdash m$ to $[t]\Iorig$ and $[t]\Irot$ agree at every Euler
grading. Their common sum in $q^{-D_m}K^0(B)[\sigma^{\pm1}][[q]]$ is
\begin{equation}\label{eq:mark}
 [t]\mathcal I_m
 =\sum_{\lambda\vdash m} Z_\lambda(q,\kappa)\,
 D\bigl(L(1-c)G_\lambda\bigr).
\end{equation}
The summand is the contribution of the indicated diagram. Its
quotients are expanded at $q=0$ before pushforward.
\end{theorem}

\begin{corollary}\label{cor:strip}
For each $\lambda$ the localized family characters agree at
$q^{-b\sum_{x\in\lambda}h_x+n}t^k$ whenever $n,k\ge0$ and
$\min(n,k)\le1$.
\end{corollary}

In \cite[Theorem~1.1]{W} the two characters at $\kappa=1$ agree
already on the common space of cycles. They are supported on the
section $m[C]$. After pushforward to $B$, Theorem~\ref{thm:first} shows that the
equality persists modulo $(1-\kappa)^2$ with the common normalization. It also extends
the genus-two, degree-two calculation of \cite[Theorem~1.2]{W},
which treated the lowest Euler grading and every $k>0$.
Theorem~\ref{thm:mark} instead retains the entire volume character
at $k=1$. Both theorems here concern classes on $B$. Lifting these
refined equalities to the common space of cycles is a further question.

The two computations of \eqref{eq:first} have different fixed loci.
On the local curve the fixed virtual sheaf has dimension zero. On
the product threefold a moment-map deformation separates the boxes
into walls whose moduli are symmetric powers of $C$. For
\eqref{eq:mark} the distinction occurs already on a disk. One
expression is a Macdonald descendant of the vertex of
\cite{DS}, the other is the sum of individual modifications at the
corners. Their equality follows by telescoping the contributions of
pairs of boundary steps. Keeping both divisors of marks in the
descent gives the stated Picard classes.

\section{Polarizations and the common expansion}\label{sec:geometry}
Let $\rho$ denote projection off the curve on any stable-pair
parameter space, and $E$ the pushforward of the sheaf to the curve.
This projection on the support is finite, being affine and proper.
Purity and flatness of the family make $E$ a vector bundle of rank
$m$. We take the relative Euler characteristic on the proper torus-fixed
locus with its virtual structure sheaf, times the inverse moving
Euler $\lambda_{-1}((T^{\rm mov})^\vee)^{-1}$ and the ordered
determinant line of \cite{NO,W}. All obstruction theories are relative
to $B$. Equivariant $K$-theoretic virtual localization is used in the
sense of \cite{Q}. We give the two lines explicitly, suppressing
constant one-dimensional vector spaces with trivial character.

\begin{lemma}\label{lem:root}
Write $\mathcal D(E)=\det R\rho_*E$ and use angular brackets for the
Deligne pairing on the curve. With $N=tL$, the original polarization,
virtual tangent and modified determinant at Euler characteristic $j$
are
\[
 P=R\rho_*[E-(1-N)\End E],\qquad
 T^{\rm vir}=P-\kappa P^\vee,\qquad
 \sigma^{j+\rk P}(\det P)^{-1}.
\]
On the rotated space the virtual tangent and modified determinant are
\[
 R\rho_*[E+cE^\vee-(1-q)(1-c/q)\End E],\qquad
 c^{k+mb/2}\mathcal D(E)^{-1}\langle\det E,L/K\rangle .
\]
The counting signs are $(-1)^j$ and $(-1)^k$, respectively.
\end{lemma}
\begin{proof}
The two coordinate Koszul resolutions give the endomorphism terms,
and the section and its Serre dual give the two linear terms.
In the first presentation curve duality pairs the two halves since
$N(\kappa K/N)=\kappa K$. The chosen square root of
$\det(T^{\rm vir})^{-1}$ is $\sigma^{\rk P}/\det P$ and the ordered
external ratio contributes $\sigma^j$ after removing the counting
variable. Riemann--Roch gives $\rk P=j+m^2b$.

In the second presentation the endomorphism term with its two
Koszul factors has trivial determinant. The determinant of the
linear terms, combined with the external ratio, has underlying
square
\[
 \frac{\mathcal D(EL)}{\mathcal D(EM)\mathcal D(E)\mathcal D(E^\vee)}.
\]
For a line $H$, the identity
$\mathcal D(EH)/\mathcal D(E)=
 (\mathcal D(H)/\mathcal D(\OO))^m\langle\det E,H\rangle$
follows from the determinant definition of the Deligne pairing and
additivity (or a splitting of $E$). Curve duality gives
$\mathcal D(E^\vee)=\mathcal D(EK)$ and
$\mathcal D(M)=\mathcal D(L)$. Substitution gives the square of
$\mathcal D(E)^{-1}\langle\det E,L/K\rangle$ up to a constant line.
Taking this ordered expression chooses the root in families. Its
character exponent is $(\chi(EM)-\chi(E^\vee))/2=k+mb/2$.
\end{proof}

\begin{lemma}\label{lem:ring}
In either presentation all coefficients belong to
$q^{-D_m}K^0(B)[\sigma^{\pm1}][[q,t]]$. Multiplication by $\sigma^{mb}$
makes the powers of $\sigma$ even.
\end{lemma}
\begin{proof}
First localize only at the antidiagonal subtorus of the internal
plane. Every geometric fixed support lies over the zero section,
since the connected torus fixes each point in a finite support
fiber. The stable-pair support morphism is proper, by the projective
bundle compactification and restriction to cycles disjoint from
infinity. Thus the fixed scheme is proper over $B$, also with its
nilpotents. Each nonzero moving weight has a unit extremal Euler
coefficient, a signed line. Its Laurent expansion can be pushed
forward coefficientwise with the residual volume equivariance still
present. On $B$ this gives Laurent polynomials in $\sigma$.
Successive localization agrees with full-torus localization. It
suffices to check the bounds after inverting $1-\sigma^r$,
$r>0$, since these are non-zero-divisors even over a ring with
additive torsion.

On the original full fixed locus the box $(i,j)$ gives the
bundle $L^{-i}M^{-j}\OO(D_{ij})$ with its coordinate character.
Thus the Euler exponent is $-b\sum h_x+\sum|D_{ij}|$.
Here $\sum h_x=\sum(1+i+j)\le m(m+1)/2$ by a linear extension of the
diagram. Reflecting a negative $t$-weight in a quotient of two Euler
classes with the same bundle changes its constant by a volume unit.
There are no negative $t$ powers. On the rotated side $k=\sum |D_{ij}|$.
To bound the $q$ exponent, forget volume weights in the \emph{ranks}
of the moving tangent. Put
$V=\sum q^{j-i}$ and $P_1=\sum |D_{ij}|q^{j-i}$, and write
$X^*=X(q^{-1})$. The change from the empty tangent has rank polynomial
\[
 F=P_1-P_1^*-(1-q)(1-q^{-1})(P_1V^*-VP_1^*).
\]
Its coefficients are antisymmetric and $\sum_{i>0}i[q^i]F=P_1'(1)$.
The reflected moving Euler has this change in its lowest exponent.
The determinant root changes it by $-P_1'(1)$, while the ordered
external lines have equal degrees. The empty tangent is $-b$ times
the Hilbert tangent, with positive exponents $h_x$, giving the bound.
The two root exponents in Lemma~\ref{lem:root} differ from $-mb$ by
even integers in the $\sigma$ variable. This proves the assertion.
\end{proof}

\section{The first variation on the local curve}\label{sec:original}
\begin{proposition}\label{prop:original}
The original index satisfies \eqref{eq:first} on every transpose orbit.
\end{proposition}
\subsection{The fixed virtual zero-cycle}


Fix a theta characteristic $S$ as reference target, $S^2=K$. For the original antidiagonal action, tensor the weight-$j$ piece of a fixed pair at $S$ by $(L/S)^j$ to obtain the full Picard family. In each weight-$j$ summand the inverse twist is $(S/L)^j$. Both operations preserve the section and coordinate maps, commute with base change, and preserve Euler characteristic since $L/S$ has degree zero on the curve. They identify the family of fixed schemes with the product over $B$. In a local frame of $L/S$, each weight-zero map of the trace-free pair complex has total twist zero. The relative Atiyah class in the pair-parameter directions holds $L/S$ fixed. Thus the fixed obstruction morphism and virtual sheaf are pulled back from the reference scheme \cite[\S2]{W}. We use the residual volume action before specialization only to perform equivariant pushforward and form the moving Euler ratios.

At the function field of $C$ a fixed pair is a finite cyclic homogeneous quotient of the two-variable polynomial algebra supported at the origin. Every weight space of that algebra is a monomial times the one-variable ring in the product of coordinates. Its homogeneous ideal in that space is generated by a power of this product, since the product acts nilpotently. Hence the generic quotient is monomial. Its diagram $\lambda$ is determined by the ranks of the weight bundles $E_j$ on $C$: on each diagonal of the diagram the occupied boxes form an initial segment. These ranks give a decomposition into open and closed loci $F_{\lambda,n}$ at Euler grading $n$. Swapping the two copies of $S$ carries $F_{\lambda,n}$ to $F_{\lambda^T,n}$.

On such a locus the fixed tangent at volume one is $P_0-P_0^\vee$ for the zero-weight part of the polarization $P$ of Lemma~\ref{lem:root}. Its virtual dimension is zero. Its ordinary virtual structure sheaf belongs integrally to the dimension-zero filtration of $G_0(F_{\lambda,n})$. We recall the integral filtration argument. For a vector bundle $p:V\to X$ of rank $r$, flat pullback carries $F_jG_0(X)$ onto $F_{j+r}G_0(V)$. Indeed Chow groups map onto the associated graded groups for the dimension filtration by taking structure sheaves, compatibly with flat pullback, and integral Chow homotopy invariance identifies $CH_j(X)$ with $CH_{j+r}(V)$. Subtract a lift in highest degree and descend in the filtration. In particular the pieces below dimension $r$ on $V$ vanish. Homotopy invariance of $G_0$ makes the inverse the zero-section Gysin map. For a global two-term bundle resolution of the fixed perfect obstruction theory (available on the quasi-projective scheme), the lifted intrinsic normal cone in the obstruction bundle of rank $r$ has dimension $r+\mathrm{vdim}$. Its structure sheaf pulled back along the zero section gives the assertion. This argument uses ordinary coherent $K$-theory after forgetting the residual torus \cite[\S3]{FG}.

Write $e_{\lambda,n}$ for this class with the counting sign, the moving spin sign and the fixed determinant line at volume one included. It is an integral sum of closed point classes. The line does not affect a point class. The sum over disconnected components retains their respective signs. Transposition preserves $e$ under transport. In fact the two polarizations differ by $V+V^\vee$ at volume one by curve Serre duality, with $V$ the push of the coordinate-weighted endomorphism bundle up to sign. This leaves the determinant of the fixed part and parity of the moving rank unchanged. The swap transports the fixed stable-pair obstruction theory itself. The degree of $e_{\lambda,n}$ is $[q^n]Z_\lambda(q)$. Here the numerical input is available \emph{partition by partition}: combining the three antidiagonal universal series of Monavari at balanced levels cancels the factor $F_\lambda$ and gives the hook product with the signed counting convention of \cite{W}. The generic diagram of each full-torus fixed component in that computation is exactly the one determined by the ranks above. This uses a numerical calculation only for the degree of an already integral zero-cycle class \cite[Theorem~1.7 and \S9]{M}.

\subsection{Positive exchanged grading}

A nonzero moving weight $t^j V$ in $P$ has spin-normalized Euler ratio $\lambda_{-\kappa^{-1} t^j}(V)/\lambda_{-t^j}(V)$, where $\lambda_z(V)=\sum_{d\ge0}z^d[\bigwedge^dV]$, after extraction of signs and scalar characters. At negative $j$ reflect the ratio to a positive power with dual coefficient and inverted volume parameter. These identities hold equivariantly on the proper fixed scheme. Every nonconstant coefficient of a normalized ratio is divisible by $1-\kappa$ by the splitting principle, and its constant term at $t=0$ is a volume unit. Divide the coefficient at $t^k$, $k>0$, by $1-\kappa$ and then specialize. Only one moving ratio contributes. A positive $j\mid k$ supplies $-\psi^{k/j} V$; a negative $-j$ with $j\mid k$ supplies $+\psi^{k/j} V^\vee$. All other factors specialize to their zeroth order, giving $e_{\lambda,n}$. This order of operations requires neither a volume derivative of a residual equivariant sheaf nor division by two. The actual graded family insertions can be restricted to closed points of the reference fixed scheme before pushforward by the dimension filtration argument above.

On the minimal quotient the plane Hilbert tangent gives the positive-half monomials $N^{a_x+1}P_2^{-\ell_x}$, where $N=tL$, $P_2=K/(tL)$ and $a_x=h_x-1-\ell_x$. The polarized tangent at the minimum is their curve push modulo hyperbolic summands, by the following comparison. The Hilbert tangent at a monomial ideal has pairs $w,NP_2/w$. Their curve pushforwards are Serre dual with opposite signs. The difference between the positive-half polarization and the Koszul polarization before pushforward is antisymmetric under $w\mapsto NP_2/w$. There is no fixed monomial in the two independent normal variables, so the difference is a sum of opposite pairs. They push to hyperbolic summands. A hyperbolic summand $A+A^\vee$ contributes zero to the positive-$t$ derivative. The resulting minimal insertion is the negative inner sum of Theorem~\ref{thm:first}.

To compute its change at an arbitrary closed fixed pair, the section and its monomials embed the minimal quotient bundle $E^{\min}$ into $E$. Monomials outside the diagram vanish generically and hence everywhere by torsion-freeness of $E$. The map on the remaining bundle is generically an isomorphism, so its kernel is zero. This also applies when a weight has rank greater than one. For a curve point $p$ let $d_j(p)$ be the length of the weight-$j$ part of the finite cokernel there. Write
\[
V(z)=\sum_j\operatorname{rk}(E_j)z^j,\qquad Q(z)=\sum_j d_j(p)z^j,\qquad f^*(z)=f(z^{-1}).
\]
Using $z=tL_p$ in the family $K$-group, the change in the Koszul polarization from this support point is
\[
F=Q-(1-z)(QV^*-VQ^*).
\]
Indeed in $K^0(C)$ finite-length duals acquire a minus sign, a point squared is zero, and multiplication by an ordinary bundle supported against a point sees its rank. The constant lines $S_p,K_p$ have class one. All family weights are precisely the powers of $L_p$. Thus the insertion change is the sum, for $j>0$ dividing $k$, of $[z^j]\bigl(F^*-F\bigr)$ times $L_p^k$. Explicitly
\[
F^*-F=Q^*-Q+(2-z-z^{-1})(QV^*-VQ^*).
\]
It changes sign on simultaneously inverting $Q,V$. This proves cancellation between distinct transposed shapes with their transported zero-cycles.

For a self-transposed shape $V=V^*$, the right hand side is $(Q^*-Q)(1-(2-z-z^{-1})V)$. Each insertion is now an integer combination of differences $d_j(p)-d_{-j}(p)$. One must justify cancellation without choosing a transposition-invariant representative of a zero-cycle. On the reference fixed scheme form the perfect insertions
\[
J_{j,k}=R\rho_*[(E_j-E^{\min}_j)L^k],
\]
using the reference weight bundles on $C$, with their torus weights omitted, and pulling $L$ from $C\times B$. At each fixed pair their restrictions are $\sum_p d_j(p)L_p^k$. The construction needs the minimal map only at closed points; the class in brackets is globally defined by a difference of bundles. Swapping the coordinates on the reference target carries $J_{j,k}$ to $J_{-j,k}$, keeping the Picard factor fixed. Their difference has zero push against the invariant class $e_{\lambda,n}$ by the projection formula and the change of variables under the swap. For distinct transposes the same global insertions can be used on transported representatives. This argument cancels actual differences in integral $K^0(B)$, including any possible additive torsion.

What remains on each orbit is therefore the minimal insertion on each shape multiplied by the degree of its fixed zero-cycle. Summation in $n$ proves the formula.

\section{The first variation on the product threefold}\label{sec:walls}
\subsection{The graded moment equations}

Retain the $q$-antidiagonal circle and decompose $E=\sum_i q^i E_i$. At the central parameter the dimension vector is $v_i=\#\{(x,y)\in\lambda:y-x=i\}$. Put $e_i=\deg E_i$ and $k=\sum e_i$. The $q$-fixed framed matrix data are the doubled linear quiver on the occupied integers with a single framing at vertex zero. There is a bundle map in each direction along each edge, a map from the trivial framing and a map back to it. A path applied to the framing is a global section. Stability means the framing generates over the function field of $C$. The moment equation at vertex $i$ can now be deformed to $a_i\,1_{E_i}$, with independent scalars of volume weight $c$. The original single-parameter deformation is the diagonal $a_i=x$.

On the central fiber the equivalence between generically cyclic framed ADHM bundles and stable pairs is \cite[Theorem~1.9]{Di}. The two plane coordinates act as the two matrices and the stable-pair section gives the framing. The return arrow vanishes for a cyclic commuting pair, and extends by density from the generic point. For the deformation retain this arrow and the full moment relation $[X,Y]+ij=a$. On a smooth parameter base use the relative obstruction theory obtained by pushing the infinitesimal gauge, arrow-and-framing, and moment complex along $C$. The transpose of the last arrow is the infinitesimal stabilizer by the trace symplectic form, with the moment character included. Consequently the first kernel vanishes and the last cokernel is supported at finitely many curve points, by generic cyclicity. The push has tangent amplitude $[0,1]$. More precisely, let $S=\AAo^2_{x,y}$ with the parameter map $a=(a_i)$ specified below, let $G=\prod_i{\rm GL}(v_i)$, and let $A$ be the affine space of graded arrows and the two framing maps for vector spaces of ranks $v_i$. Write $\mathfrak g^*$ for the moment target, using the trace pairing. Over $S$ take
\[
 \mathfrak X=\big[(A\times S)\mathop{\times}^{\mathbf R}_{\mathfrak g^*} \{0\}\,/G\big],
 \qquad (A\times S)\longrightarrow \mathfrak g^*,\quad (z,s)\longmapsto\mu(z)-a(s).
\]
The relative derived moduli is the open locus of generically cyclic maps of the specified bundle degrees in $\underline{\rm Map}_S(C\times S,\mathfrak X)$. Indeed an $S$-test object gives bundles on its product with $C$, their arrows and a derived solution of the moment relation. For any derived $S'\to S$, this same description on each $S'$-test object identifies the base change with maps to $\mathfrak X\times_S^{\mathbf R}S'$. The stability condition and degrees restrict on classical fibers. Thus the asserted fiber squares are derived Cartesian even for nontransverse parameters. The relative tangent is the curve pushforward of the three-term gauge, arrow and moment complex above. On the central fiber its fixed virtual tangent is the zero-$q$ part of
\[
R\rho_*[E+cE^\vee-(1-q)(1-c/q)\operatorname{End} E],
\]
and its fixed virtual sheaf is the induced fixed stable-pair sheaf. This follows by taking equivariant weight zero in the equivariant Koszul complex of the pair: its gauge term is $\End E$, its arrow and framing term is $(q+c/q)\End E+E+cE^\vee$, and its moment term is $c\End E$. Alternatively the intrinsic Fulton class formula \cite[Theorem~4.2]{Th} identifies the virtual sheaves on the identical fixed scheme with identical virtual cotangent classes. The nonzero-$q$ part of this expression defines moving Euler insertions on the entire parameter family, even at parameters specified only on the fixed quiver. The ordered determinant also extends by determinants of cohomology and the external line pairing. At $c=1$ it is $\mathcal D(E)^{-1}\langle\det E,LK^{-1}\rangle$ up to constant lines, $\mathcal D=\det R\rho_*$. The counting sign is $(-1)^k$. These statements require no deformation of an ungraded stable-pair space in the independent vertex parameters.

Choose vertex parameters $a_i=x+y\eta_i$ with $x,y$ both of character $c$. Choose $\eta_i$ so that all averages on distinct consecutive intervals are distinct. For $y\ne0$ each value of $x$ lies on at most one wall $\sum_{a\le i<d}a_i=0$. For $y=0$ and $x\ne0$ all these sums are nonzero. The same construction works with only one vertex.

\subsection{Reflections and properness}

Introduce spectral parameters $s_i$ with $a_i=s_i-s_{i+1}$ and bits
\[
b_i={\bf1}_{i\le0}-v_i+v_{i-1}.
\]
Here the finite range of indices runs from the first occupied vertex to one past the last, with outside dimensions zero. The $b_i$ are zero or one, from the diagonal counts of a Young diagram. At a vertex with $a_i\ne0$ reflection exchanges $s_i,s_{i+1}$ and $b_i,b_{i+1}$. On bundles it replaces $E_i$ by the kernel of $A:H\to E_i$, where $H=E_{i-1}\oplus E_{i+1}\oplus\mathcal O^{\mathbf1_{i=0}}$. For the signed outbound map $B_i:E_i\to H$ the moment equation is $AB_i=a_i1$. Thus $H=B_iE_i\oplus\ker A$ in families. Replace the two maps by inclusion and $-a_i$ times projection for the complementary summand, so their product on $H$ changes from $B_i A$ to $B_i A-a_i1_H$. The neighboring parameters become $a_{i\pm1}+a_i$ and the vertex parameter becomes $-a_i$. Signs on the incident arrows implement the chosen moment signs. This construction is involutive up to canonical isomorphism. It preserves generic cyclicity: it gives an equivalence on representations with the framing vertex retained and on their subquotients, by the same split kernel construction. One may impose the moment relation at the framing vertex as well: its scalar there is forced by the sum of traces, since the framing has rank one. The construction bypasses that vertex. In particular the degree charges $e_{i-1}-e_i$ are interchanged at the same two positions. This is the nonzero-parameter reflection for a deformed preprojective algebra, applied to vector bundles.

When no spectral values coincide, sort the bits with all ones first. All gauge ranks then vanish, and the unique generically cyclic representation has zero degree. Suppose exactly $s_a=s_d$, $a<d$. All swaps are allowed except interchanging these two labels. Unless $b_a=0,b_d=1$ one can still sort to the vacuum. In the exceptional case put all other ones first, then the two equal labels in their given order, then all other zeros. The only remaining gauge vertex is the framed one and it has rank one. At the wall the one equation is $ij=0$. In a transverse family it is $ij=a'_0$, where $a'_0=s_a-s_d$ up to the fixed sign convention. The gauge line at positive degree is $\mathcal O(D)$, its section gives the effective divisor $D$, and the return map is zero. Transporting the degree charges back gives
\[
e_i=d_0\,\mathbf1_{a\le i<d},\qquad d_0=|D|>0,\qquad k=(d-a)d_0 .
\]
On this component the same inverse reflections give, in ordinary $K$-theory over the full divisor family,
\[
[E_i]=v_i+\mathbf1_{a\le i<d}([\mathcal O(D)]-1).
\]
Indeed the reflection relation for classes is $[E_i']=[H]-[E_i]$. Its coefficient of $\mathcal O(D)-1$ transforms by the same simple reflection as the degree. All calculations at $y\ne0$ here are after forgetting the volume action.

\begin{lemma}\label{lem:proper}
At fixed ranks and degrees the graded moment family is proper over the parameter plane. At positive total degree it is proper over the $y$-line.
\end{lemma}
\begin{proof} The finite-dimensional, cyclic, \emph{constant} framed representations with dimension $v$ form an \'etale family over the $(x,y)$-plane. Cyclicity implies trivial infinitesimal stabilizer and surjective moment differential, and the relative tangent dimension is $2(v_0-\sum v_i^2+\sum v_i v_{i+1})=0$. The dimension identity follows from $b_i^2=b_i$ and summation, or from the absence of weight zero in the monomial Hilbert tangent. The stable free gauge quotient carries its tautological bundles. There is exactly one geometric point over every parameter: off intersections use the reflections just given and the unique cyclic framed rank-one or vacuum representation. At the origin a cyclic framed ADHM point has zero return map and is a homogeneous colength-$m$ ideal; its monomial scheme is determined by $v$. Thus this \'etale family is an isomorphism to the parameter plane (an \'etale universally injective map is an open immersion). Denote its cyclic bundles and arrows by $T_i$. Their constant pullbacks are the minimal bundles on the curve.

For a generically cyclic bundle representation over a field, words on the framing embed the minimal $T_i\otimes\mathcal O_C$ into $E_i$, generically isomorphically. Indeed the generic representation is the unique stable point, canonically identified by its framing. Relations among words extend across missing curve points by torsion-freeness. The dual inclusions have finite quotient of length $e_i$. This description holds over nonreduced bases. Locally choose a basis of words for the constant cyclic bundle. The corresponding sections of $E_i$ satisfy the same relations on the locus of generation, by uniqueness of the framed cyclic representation there. This open is schematically dense relatively: its complement is fiberwise finite, and a local determinant that is regular on each curve fiber is regular on the flat bundle. Relations therefore extend on the whole family. The dual word map gives a length-$e_i$ flat quotient, by the fiberwise injection criterion. Conversely kernels of such quotients are bundles, and the incidence conditions that the dual arrows preserve them and the dual return section lies in the kernel are closed. Dualizing recovers the framed maps with their relations; the word sections ensure generic generation. Equivalently the valuative criterion reads as follows: over a discrete valuation ring extend the colored dual quotients by the proper relative Quot schemes of the $T_i^\vee$. The extended kernels are bundles by flatness and torsion-freeness on curve fibers. The dual arrows preserve them by closedness of composition to the finite flat quotients. The dual of the minimal return map factors through the extended kernel by the same argument. The section extends from the minimal bundle automatically. The moment equations hold by density and flatness. Each fiber generates generically as it contains the minimal word sections. Uniqueness follows from the forced word maps in any extension and separatedness of Quot. Finite type and separatedness of the stable framed moduli can also be seen on the bounded word charts: the ranks and the degrees are fixed, the generically independent word sections put each dual bundle in a finite length Quot of a trivial bundle, and the arrow and relation spaces over the resulting bounded families are of finite type. The moduli is proper over the parameter plane at fixed degrees. For $k>0$ its support is on a finite union of lines finite over $y$. The structural morphism to the $y$-line is therefore proper as well, including nilpotent thickenings. One may use a common power of the product of the wall equations to cut out the scheme-theoretic image.

\end{proof}

\subsection{Taking the first variation by a transverse line}

\begin{lemma}\label{lem:basechange}
For $k>0$, dividing the central index by $\delta$ and setting $c=1$
gives the ordinary-volume index over the line $y=1$, with the
restricted moving factors and ordered root.
\end{lemma}
\begin{proof}
Let $\mathfrak M$ denote the relative derived moment family. The two
successive base changes, with $a$ equal to $0$ or $1$, are
\[
\begin{CD}
 \mathfrak M_{y=a}@>>> \mathfrak M\\
 @VVV @VVV\\
 \AAo^1_x @>{y=a}>> \AAo^2_{x,y}
\end{CD}
\qquad
\begin{CD}
 \mathfrak M_{x=y=0}@>>>\mathfrak M_{y=0}\\
 @VVV @VVV\\
 \{0\}@>{x=0}>>\AAo^1_x .
\end{CD}
\]
The two squares are derived Cartesian by the functor of maps just described.
Their restricted obstruction theories are those of the indicated fibers:
the curve is fixed and proper and the gauge--moment complex has
perfect pushforward compatible with base change.
The virtual pullbacks thus commute with the proper pushforwards and
with tensoring by the perfect insertions \cite[Theorem~2.6]{Q}.
Push the fixed virtual family on the plane to the $y$-line, including the moving Euler and root insertions coefficientwise in $q$ and the constant Picard base $B$. At $y=0$ this gives the virtual index of the total diagonal $x$-family by virtual base change and proper pushforward. The latter is supported over $x=0$ for $k>0$. Restriction along $x=0$ multiplies its pushed index by $1-c^{-1}$: filter each sheaf by the nilpotent parameter and resolve the origin of the character-$c$ line by its two-term Koszul complex. The insertions are restrictions of perfect classes, with the moving Euler inverse expanded at $q=0$. Consequently the central family index divided by $\delta$ at $c=1$ is the ordinary-volume push from this total $x$-family. There is no residual factor from the number of vertices. Using a single transverse line at this stage fixes the normal character, independent of the length of a wall.

After division, forget the volume action. The fibers $y=0$ and $y=1$ of the pushed class agree by homotopy invariance on $B\times\mathbb A^1_y$ and virtual base change. Thus one may calculate on the affine line $y=1$, where the surviving wall points are distinct. Near each contributing point, the ordinary moduli on the $x$-line is the smooth scheme $C^{(d_0)}$. The reflections above are isomorphisms over the open set where their parameters are invertible. In the remaining rank-one problem, $H^0(C,\mathcal O(-D))=0$ in all fibers, and the degree is positive. Relative cohomology and base change force the return map to vanish also in families. Hence $a'_0=0$ scheme-theoretically. This uses the simple zero of the transverse linear parameter.

\end{proof}

\subsection{Contribution of one wall}

Let $p_D:D\to C^{(d_0)}$ be the universal divisor projection. Write $V=\sum v_iq^i$, $W=\sum_{a\le i<d}q^i$, $h=d-a$, $U=p_{D*}\mathcal O_D$, $T=p_{D*}\mathcal O_D(D)$, $J=\det U$, and $F^*=F(q^{-1})$ for a Laurent polynomial. The quotient on the wall has class $V+W(\mathcal O(D)-1)$. Inserting this into the framed tangent at $c=1$ gives the increment $TH-UH^*$ with
\[
H=W[1-(1-q)(1-q^{-1})(V-W)^*]=\sum_j h_jq^j .
\]
The empty tangent is $-b$ times the Hilbert tangent, which has no weight zero. One has $h_0=1$, $H(1)=h$ and $H'(1)=W'(1)$. The fixed tangent over the transverse line is therefore $T-U+1$ in $K$-theory. Since the underlying scheme is smooth with tangent $T$, its obstruction bundle has class $U-1$. Its virtual sheaf is $\lambda_{-1}((U/1)^\vee)$. This determines the fixed virtual sheaf without asking a reflection to identify the moving stable-pair obstruction theories.

Here are the coefficient and hook identities explicitly. Put $\mu_i=v_i-\mathbf1_{a\le i<d}$. Summing the second differences of $\mu_i$ gives
\[
h_0=\mathbf1_{a\le0<d}-(\mu_a-\mu_{a-1})+(\mu_d-\mu_{d-1})=1 .
\]
The sum of the coefficients of the positive powers of $q$ in $H$ equals
\[
p=\sum_{j=a}^{d-1}(\mathbf1_{j>0}+\mu_j-\mu_{j-1})
=\#\{i:a<i<d,\ b_i=0\}.
\]
The inversion pairs $b_a=0,b_d=1$ are in bijection with boxes of $\lambda$. In fact the zero corresponding to row $y$ is at $a=y-\lambda_y+1$ and the one corresponding to column $x$ is at $d=\lambda'_x-x$, with rows and columns starting at zero and lengths denoted by $\lambda_y,\lambda'_x$. These positions are the vertical and horizontal steps of the partition boundary, or follow by substituting the diagonal counts in the bit formula. The inversion condition is precisely $(x,y)\in\lambda$. It gives $h=d-a=h_{(x,y)}$, and the intervening zeros are the rows below $y$ reaching column $x$, so $p=\ell_{(x,y)}$.

Write $N_D(F)=\det p_{D*}(F|_D)/\det U$ for the norm of a line. Lemma~\ref{lem:root} gives the relative ordered root at $c=1$, which is $(\det W)^{-d_0}J^h N_D(L^h)$ by the two divisor sequences and the Deligne pairing. Reflect the dual positive-weight Euler denominators of $TH-UH^*$ about $q=0$. The character $(\det W)^{-d_0}$ cancels by $H'(1)=W'(1)$. Since $\det T=J^{-1}N_D(K)^{-1}$, their determinants leave $J\,N_D(L^hK^{-p})$ after the root. Their sign, the counting sign $(-1)^{hd_0}$, and conversion of $J\lambda_{-1}((U/1)^\vee)$ together give $-1$. Indeed the positive moving rank increment is $d_0(2p-h+1)$ and the last sign is $(-1)^{d_0-1}$. The empty quotient contributes $Z_\lambda$. The remaining normalized symmetric powers at $q^i$, $i>0$, have inputs
\[
h_i(T-U^\vee)+h_{-i}(T^\vee-U).
\]
These have rank zero and are dualized by transposing the marked box and shape. Let $F_\gamma$ denote the gamma filtration on the ordinary $K$-ring of $C^{(d_0)}$. A rank-zero class and its dual have opposite images in $F_\gamma^1/F_\gamma^2$, detected by their determinants. The positive-degree coefficients of the normalized symmetric algebra have rank zero, and dualizing its inputs dualizes its coefficients. Thus each nonconstant correction cancels with its transpose modulo $F_\gamma^2$. The norm lines in the two terms have the same $N_D(L^h)$; the discrepancy in powers of $N_D(K)$ does not alter the cancellation in the first filtration quotient. At an invariant marked box $h_i=h_{-i}$, and each input already has zero rank and trivial determinant. Multiplication by $\lambda_{-1}(U/1)$ in filtration $d_0-1$ kills all remaining corrections on the smooth $d_0$-dimensional symmetric power. Indeed $F_\gamma^{d_0+1}$ maps to zero by the dimension filtration on this smooth $d_0$-fold. This calculation precedes the family norm insertion and requires no division by an integer. For any line $F$ on $C\times B$ one has
\[
R\pi_*\bigl(\lambda_{-1}(U/1)N_D(F)\bigr)=\psi^{d_0}D(F).
\]
Indeed let $\nu:C^{d_0}\to C^{(d_0)}$ be the ordered cover, and $V_0$ the standard representation of $\mathfrak S_{d_0}$, the permutation representation modulo constants. The covariants of the permutation representation in $\nu_*\OO$ give $U$. Their exterior powers agree with the covariants of the exterior powers: locally this is Solomon's theorem \cite{So} for the product of symmetric groups on the coordinates at each support point. Split off the invariant line. The norm pulls back to $F^{\boxtimes d_0}$, so relative K\"unneth computes the push as the alternating sum of the Schur functors of $D(F)$ for the representations $\bigwedge^j V_0$. Their characters correspond to hooks $(d_0-j,1^j)$ and the integral symmetric-function identity $\sum(-1)^j s_{(d_0-j,1^j)}=p_{d_0}$ gives the Adams operation, with graded tensor permutations for $D(F)$.

Summing the surviving inversions for $k=hd_0$ proves the formula.

\subsection{First order in the constant exchanged grading}

The two character expansions also agree at $t^0$ on the first integral infinitesimal neighbourhood of $\sigma=1$, $\sigma^2=\kappa$. For a transpose orbit let $z_n=[q^n]\sum_{\lambda\in\mathcal O}Z_\lambda(q)$. On either side the coefficient of $q^nt^0$ there is
\[
z_n+nz_n(\sigma-1)\pmod{(\sigma-1)^2}.
\]
In particular no half-integral parameter derivative is needed.

For the rotated empty quotient the exact contribution of one shape is
\[
(-1)^{mb}q^{-b\sum h_x}\kappa^{-b\sum(\ell_x+1/2)}
\prod_x(1-\kappa^{\ell_x}q^{h_x})^b(1-\kappa^{\ell_x+1}q^{h_x})^b .
\]
This follows from the constant virtual tangent $-b$ times the Hilbert tangent and the ordered root $c^{mb/2}$. Denote it by $Z_\lambda(q,\kappa)$. Interchanging arms and legs shows $Z_\lambda(q,\kappa)=Z_{\lambda^T}(\kappa q,\kappa^{-1})$ on the cover, with inversion of $\sigma$ understood. Thus the orbit sum at $q^n$ is a Laurent polynomial $H_n(\sigma)$ satisfying $H_n(\sigma)=\sigma^{2n}H_n(\sigma^{-1})$.

For the original calculation at $t^0$ the antidiagonal fixed-scheme product identification at the reference spin line extends equivariantly for the remaining torus. All fixed maps and the fixed orientation carry total power zero of the varying Picard line. Reflection of the nonzero moving Euler factors leaves only their scalar powers of $\sigma$ in the $t^0$ coefficient. Hence this coefficient is an integer Laurent polynomial times $[\mathcal O_B]$, and can be computed on the reference target. At fixed $n$ remove its external factor $\sigma^n$ and consider the symmetrized virtual index at variable $t$. It is invariant under simultaneous inversion of both torus characters. Indeed on a proper full-torus fixed component with fixed virtual dimension $d$, apply weak virtual Serre duality to the determinant-root and inverse moving Euler insertion. Dualizing the moving Euler contributes its determinant and the sign $(-1)^{\mathrm{rk}N^{\rm vir}}$, by the line identity $1-z=-z(1-z^{-1})$. The determinant cancels against the square of the total root and the fixed virtual canonical bundle. The resulting total sign is $(-1)^{d+\mathrm{rk}N^{\rm vir}}=1$ by the virtual dimension of the local Calabi--Yau threefold theory. Ordinary virtual Serre duality on the fixed scheme applies coefficient by coefficient with the torus weights recorded as scalars; equivalently one can clear the Euler denominators \cite[Proposition~3.10]{FG}.

The target also admits the coordinate swap because both normal lines at the reference point are the same spin line. The ordered polarization changes by a hyperbolic summand and the external factor does not depend on which normal coordinate was chosen. On the orbit sum the swap takes $t$ to $\kappa/t$. Composing it with inversion takes $(t,\sigma)$ to $(t/\kappa,\sigma^{-1})$, preserving the chamber at $t=0$. Restoring the external factor and extracting the constant term therefore gives the same reciprocity $H_n(\sigma)=\sigma^{2n}H_n(\sigma^{-1})$ on the original side. On both sides $H_n(1)=z_n$. Differentiating at one yields $2H'_n(1)=2nz_n$ in the integers, which gives the stated coefficient before multiplying by the unit class of $B$.

In fact, after extracting the common factor $\sigma^{-mb}$ the two full characters are Laurent-polynomial in $\kappa$ coefficientwise and agree modulo $(1-\kappa)^2$ over $K^0(B)$ at all gradings. For the original root the exponent of $\sigma$ is $n+\operatorname{rk}P=2n+m^2b$ up to even powers. For the rotated root its half-character exponent comes from $c^{k+mb/2}$, since $\chi(EM)=k$ and $\chi(E^\vee)=-k-mb$. All other character weights come from the original torus. Thus multiplication by $\sigma^{mb}$ removes the determinant cover on both sides. At $k>0$ the proof computed the quotient by $1-\kappa$ itself, so equality of that quotient at one implies the required congruence directly in the Laurent ring. At $k=0$ the coefficients are scalar integers, and the equality of values and first derivatives in $\sigma$ gives equality of the first derivatives of the corresponding polynomials in $\kappa$ after removing the common factor. Here cancellation of two is in $\mathbb Z$ only. This yields the all-grading first-order comparison without any assumption about two-torsion of the Picard $K$-group.

\section{The stable-pair vertex on a disk}\label{sec:disk}
We work temporarily on an affine line with cotangent character $p$.
The parameters $p,c$ are independent and $s=p/c$. The Poincar\'e
weight is included in $y$; on the curve $x=q$ and $y=tL$ in a
local frame. The coefficient of $x^ny^k$ measures excess above
the minimal Euler grading of this particular diagram. We give both
the unsigned boundary measure and the determinant factors that
relate it to the geometric vertex.
\subsection{The polarized character}


A reverse plane partition is a weakly increasing nonnegative array $d_{ij}$ on the boxes $(i,j)$, with $(0,0)$ at the root. Write $|d|=\sum d_{ij}$. Set
\[
 E_d=\sum_{(i,j)\in\lambda} y^{j-i}s^{-j}p^{-d_{ij}},
 \qquad H_y(E)=\frac{E-(1-y)EE^\vee}{1-p},
 \qquad \Delta H_d=H_y(E_d)-H_y(E_0).
\]
A dual inverts all characters. Write $(z;p)_d=\prod_{i=0}^{d-1}(1-zp^i)$ and use the reciprocal shifted product for negative $d$. Put $R_c(w)=(1-cw)/(1-w)$ and extend by products to finite virtual characters. The full formula is
\[
 A_{\lambda;c,s}(x,y)=
 \prod_{\square\in\lambda} F_c(y^{h_\square}s^{-l_\square})
 \sum_{d} x^{|d|}c^{-|d|} R_c(\Delta H_d),
 \qquad F_c(z)=\sum_{j\ge0}\frac{(c;p)_j}{(p;p)_j}z^j .
\]
The factors in $R_c$ are read at $y=0$ after reflecting negative powers. The finite difference $\Delta H_d$ is a Laurent polynomial. Its rank is $|d|$, by the divided difference of the framing term; the rank contribution from the endomorphism differences vanishes.

The classification of fully torus-fixed local pairs gives these arrays. The sheaf is a weight line with divisor $d_{ij}[0]$ at each box, and the multiplication arrows are their canonical inclusions. In particular the inequalities at a box with two predecessors apply simultaneously. For the one-leg correspondence and virtual character see \cite[Theorem~6.1 and \S6.2]{MR}. In our variables the original tangent is the polarized character $H-c^{-1}H^\vee$. For the rotated vertex one substitutes the plane weights $x,c/x$ instead, and obtains $H-s^{-1}H^\vee$ since
\[
 (1-x)+c(1-x^{-1})=(1-x)(1-c/x),\qquad
 (1-p^{-1})^{-1}=-p/(1-p).
\]
The relative symmetrized Euler, counting sign and external determinant give $c^{-|d|}R_c(\Delta H_d)$ originally, and $s^{-|d|}R_s(\Delta H_d)$ with switched inputs on the rotated disk. Indeed the internal root gives $(\det\Delta H_d)^{-1}c^{-|d|/2}$ originally. Inverting the dual Euler supplies $(-1)^{|d|}\det\Delta H_d$ and the Euler ratio. The external determinant supplies $x^{|d|}c^{-|d|/2}$ and the counting sign cancels the previous sign. For the rotated disk the external line weights are $y,s/y$ and the same calculation gives $y^{|d|} s^{-|d|/2}$ for the external ratio. The relative determinants of the finite framing change cancel in that ratio. The empty products come from the positive half of the Hilbert tangent. These statements concern the disk characters after the empty scalar on the projective curve is removed.

There are no invariant Euler factors in this prescription. For a framing contribution to the coefficient of $y^0c^0$ only the root matters, and its nonempty difference has exponents $p^{-d},\ldots,p^{-1}$. In the unshifted endomorphism terms $y^0c^0$ comes from identical boxes and has zero height difference. In the $y$-shifted endomorphism term it comes from a horizontal adjacent pair, with shift parent height minus child height at most zero, giving only strictly negative powers. Thus $1$ does not occur. For a term with $c^{-1}$ the power of $s$ must be $+1$, so the raw divided difference must contain $p^{-1}$ to give $p^0$. There is no such framing term. In the unshifted endomorphisms the two boxes are a southeast diagonal step apart, and in the $y$-shifted term they are a vertical step apart. Both height differences are nonnegative, and the associated divided differences contain powers at least zero. This rules out $c^{-1}$ at axial exponent zero. In particular both $R_c$ and its exchanged counterpart can be evaluated for generic independent $p,c$. The argument uses the order on all boxes and applies to a repeated content.

\begin{lemma}\label{lem:boundary}
The two boundary series are
\[
 A_{\lambda;c,s}(x,0)=\prod_{z\in\lambda}F_s(x^{h_z}c^{-\ell_z}),
 \qquad A_{\lambda;c,s}(0,y)=\prod_{z\in\lambda}F_c(y^{h_z}s^{-\ell_z}).
\]
Moreover the summands on the first boundary are the interlacing
partition weights of \cite[Propositions~6--7]{DS} at the shifted
parameters specified below.
\end{lemma}
\begin{proof}
The second series is the empty height array. For the first use
the following specialization of \cite[Theorem~1]{DS}.
Write a box as $u=(i,j)$ with $i$ increasing to the right and $j$ increasing downwards, starting at zero. Put $a_u=i-j$. Let $d_u$ be its height, increasing weakly rightwards and downwards. Put
\[
 B(z;D)=\frac{(sz;p)_D}{(pz;p)_D},
 \qquad r_u=y^{a_u}s^j ,
\]
with the usual reciprocal convention for a negative lower index. For a monomial contribution of weight $w$ and height difference $D$ in the polarization,
\[
 R_c\!\left(w\frac{p^{-D}-1}{1-p}\right)=c^D B(w^{-1};D).
\]
This follows by reversing the $D$ factors for $D>0$ and taking reciprocals at shifted weights for $D<0$. The powers of $c$ cancel the external $c^{-|d|}$ because the total rank is $|d|$. The three factors left for a fixed array are
\[
 \prod_u B(r_u;d_u)
 \prod_{u,v}
 \frac{B(y^{-1}r_u/r_v;d_u-d_v)}
      {B(r_u/r_v;d_u-d_v)} .
\]


At $y=0$, a factor with a positive $y$ exponent tends to one, and a factor with a negative exponent tends to $(s/p)^D=c^{-D}$. The framing therefore contributes $c^{-\sum_{a_u<0}d_u}$ from the latter group. For the ordered pairs, the numerator and denominator have negative powers simultaneously if $a_u<a_v$. The extra numerator terms with $a_u=a_v$ contribute total exponent zero, since $\sum_{a_u=a_v}(d_u-d_v)=0$. The factors with zero exponent are the framing at content zero, the endomorphisms at each content and the arrows from content $a$ to $a+1$. Their roots at content $a$ are $s^j$ with
\[
 j=\max(0,-a),\ldots,\max(0,-a)+v_a-1,\qquad
 v_a=\#\{u:a_u=a\}.
\]
These are precisely the three groups of Pochhammer factors in equation (19) of the cited hook-product paper, with $(q,\hbar)=(p,s)$ and framing parameter one. The weak height inequalities are their interlacing conditions (Proposition 7 and its proof). Thus the remaining sum equals their series $Z$ at
\[
 z_a=x\,c^{-\mathbf 1_{a<0}} .
\]
This reduction works with repeated contents. It retains the endomorphism factors between all distinct roots on a common diagonal.

In \cite[Theorem~1]{DS} each box contributes $F_s(\prod_{a=L}^{U}\widehat z_a)$, where $\widehat z_a=(s/p)^{v_{a-1}-v_a+\mathbf1_{a=0}}z_a$. The contents along the hook at $(i,j)$ fill the interval $[L,U]=[i-j-\ell,i-j+\operatorname{arm}]$ once each. Here the multiplicity diagram is indexed with positive contents to the right. Equivalently, the two endpoints record the last box in the column and the last box in the row. Reflecting the drawing interchanges the two descriptions of hooks and preserves the set of these intervals. The exponent of $c^{-1}$ in their product is
\[
 v_{L-1}-v_U+\#\{a\in[L,U]:a\le0\}.
\]
The last row on diagonal $U$ has index $j$, and the boxes on diagonal $L-1$ run from row $\max(0,1-L)$ through $j+\ell$, with the empty interval permitted. Consequently
\[
 v_U=j+1-\max(0,-U),\qquad
 v_{L-1}=j+\ell+1-\max(0,1-L).
\]
Substitution gives exactly $\ell$. There are $h$ powers of $x$. This proves the first formula over $\mathbb Q(p,c)[[x]]$. The analytic parameter range in the hook-product theorem suffices, since each Taylor coefficient is rational in the independent parameters.
\end{proof}
\subsection{The coefficient of the first positive weight}
Write $X_i$ for the alphabet $\{s^jp^{d_{i+j,j}}\}$ on diagonal $i$
and $V=\sum_i y^i X_i$. A star on a character inverts all its weights.
Set $P(V)=V-(1-y^{-1})VV^*$ and
\[
 T(V)=sP(V)+P(V)^*=sV+V^*-(1-y)(1-s/y)VV^* .
\]
Relative to the height measure at $y=0$, the positive-weight part
of the insertion in $A$ is the symmetric algebra on
$(1-c)T(V)_{>0}/(1-p)$. To see this, put
$G(z)=(sz;p)_\infty/(pz;p)_\infty$. A factor $B(z;D)$ of positive
$y$ degree is $G(z)/G(zp^D)$, and a factor of negative degree is
$c^{-D}F_c(1/(zp^D))/F_c(1/z)$. The monomials and the degree-zero
Pochhammer factors give the boundary measure of Lemma~\ref{lem:boundary}
term by term. The other factors give the symmetric algebra on the
positive weights of $(1-c)(T(V)-T(V_{\rm empty}))/(1-p)$.
At empty heights these missing weights of $T$ are
$\sum_z y^{h_z}s^{-\ell_z}$ by the arm--leg tangent formula, precisely
the prefactor in $A$. In particular the linear insertion is
$(1-c)[y]T(V)/(1-p)$.

\section{A Macdonald current at one exchanged mark}\label{sec:currents}
In this section put $x=u$, $v=c/u$, and use $R(z)=(1-sz)/(1-z)$
with $p$ still independent. We shall prove
\begin{equation}\label{eq:local}
 \frac{[y]A_{\lambda;c,s}(u,y)}{A_{\lambda;c,s}(u,0)}
 =\frac{[y]A_{\lambda;s,c}(y,u)}{A_{\lambda;s,c}(0,u)}
 =\frac{1-c}{1-p}G_\lambda .
\end{equation}
Lemma~\ref{lem:boundary} identifies the denominators.
\subsection{The single modification on the exchanged side}


At excess one a reverse plane partition is supported on a removable box \(r\). Put \(E=\sum_{h\in\lambda}d_h^{-1}\) and \(W=d_r^{-1}\). Replacing \(W\) by \(W/p\) in the polarized character changes it by
\[
 \Delta H=p^{-1}W+(1-u)\big((E-W)/W-p^{-1}W(E-W)^*\big).
\]
The relative oriented factor is \(s^{-1}R(\Delta H)\), where a sum of weights is evaluated multiplicatively. Using \(R(z)R(1/(sz))=s\), a box \(h\ne r\) contributes \(R(w)R(cw)/(R(uw)R(vw))\), \(w=d_r/d_h\), and the framing contributes \(1/R(c d_r)\). To take the boundary of this product, use the Laurent identity
\[
 (1-u)(1-v)E=c-c\sum_{a\in\mathcal{A}}d_a^{-1}
                       +\sum_{h\in\mathcal{R}}d_h^{-1}.
\]
First form the product with independent numerator variable \(z\) in place of \(d_r\). Remove the entire \(r\)-box factor before setting \(z=d_r\); the two copies of \(R(z/d_r)\) cancel. The answer is the \(r\)-summand of \eqref{eq:corners} divided by \(R(c)=(1-p)/(1-c)\). This proves the second equality in \eqref{eq:local} directly from the oriented fixed-pair polarization.

\subsection{The insertion on the original side}

Let \(X_i\) be the sum of roots \(s^j p^{d_{i+j,j}}\) on content \(i\) of a height array \(d\); a star inverts its roots. Denote by brackets the expectation in the height sum at \(y=0\), divided by its hook product. Extracting the linear coefficient of the insertion of Section~\ref{sec:disk} gives \((1-c)/(1-p)\) times \(\langle O\rangle\), where
\begin{equation}
O=sX_1+X_{-1}^*
     +\sum_i (X_i-X_{i+1})(X_i^*-sX_{i-1}^*).
 \label{eq:obs}
\end{equation}
In fact \eqref{eq:obs} is the coefficient of \(y\) in \(sV+V^*-(1-y)(1-s/y)VV^*\), \(V=\sum_i y^i X_i\).

Put \(J_i=\max(0,-i)\), \(L_i=J_i+\#\{\text{boxes of content }i\}\). The southeast boundary from \((0,\operatorname{height}\lambda)\) to \((\operatorname{width}\lambda,0)\) has vertices \((i+L_i,L_i)\) with monomials \(Z_i=u^i c^{L_i}\). A horizontal step is denoted by \(-\), a downward one by \(+\). Thus \(Z_{i+1}/Z_i=u\) at a minus step and \(1/v\) at a plus step. A removable box is precisely a minus followed by a plus. The slice partition \(\mu_i\) lists the heights on content \(i\) in decreasing order.

We use the Macdonald inner product with parameters \((p,s)\). In the power sums \(b_n\), write \(\partial_n=\partial/\partial b_n\). The right vacuum is the constant polynomial $1$ and the left vacuum takes the constant in all $b_n$. In increasing step order the normalized boundary measure is the vacuum matrix element of
\[
 \Gamma_-(Z_i)=\exp\!\sum_{n>0} Z_i^{-n}\partial_n,\qquad
 \Gamma_+(Z_i)=\exp\!\sum_{n>0}\frac{1-s^n}{1-p^n} b_n Z_i^n/n ,
\]
using the indicated operator at each step and dividing by the vacuum expectation. An insertion on slice \(i\) acts immediately before step \(i\). This is the \emph{termwise} interlacing-partition expansion of Dinkins--Smirnov, Propositions 6--7 and \S{}3.7. Indeed their shifted node parameter is \(\widehat z_i=u c^{L_i-L_{i-1}}=Z_i/Z_{i-1}\), and a minus to its left and plus to its right give the Cauchy factor \(F_s(Z_{\rm right}/Z_{\rm left})\). The common scale of the \(Z_i\) cancels \cite[\S3]{DS}.

Define the diagonal eigenvalues \(P_i,N_i\) on slice \(i\) by
\[
 P_i=(s-1)X_i+s^{J_i},\qquad N_i=-(s-1)X_i^*/s+s^{-J_i}.
\]
Write \(\Delta P_i=P_i-P_{i+1}\), \(\nabla N_i=N_i-sN_{i-1}\). Then
\begin{equation}
O=-\frac{s}{(s-1)^2}\sum_i \Delta P_i\,\nabla N_i .
 \label{eq:diff}
\end{equation}
Indeed the difference of \(s^{J_i}\) is \((s-1)s^{-i-1}\) for \(i<0\) and zero otherwise. The backward difference of \(s^{-J_i}\) is \(1-s\) for \(i>0\) and zero otherwise. Their product vanishes. The two cross terms telescope to the framing terms of \eqref{eq:obs}. Only $-\operatorname{height}(\lambda)<i<\operatorname{width}(\lambda)$ is needed. Outside the path the slice is the vacuum. At its left endpoint $\nabla N_{-\operatorname{height}(\lambda)}=0$, and at its right endpoint $\Delta P_{\operatorname{width}(\lambda)}=0$.

The normalized first Macdonald operator is the constant term in \(z\) of
\[
 \eta(z)=\exp\!\sum_{n>0}(1-s^{-n})b_n z^n/n\
          \exp\!\sum_{n>0}(p^n-1)z^{-n}\partial_n.
\]
On \(\mu\) its eigenvalue is \(1+(s-1)\sum_{j\ge1}(p^{\mu_j}-1)s^{-j}\). Let \(\xi(w)\) be the same current with \(p,s\) inverted. Both are diagonal in the same monic Macdonald basis: inversion rescales the inner product on homogeneous degree \(d\) by \((s/p)^d\), preserving triangularity and orthogonality. The two zero modes, multiplied by \(s^{L_i}\) and \(s^{-L_i}\), respectively, give \(P_i,N_i\) since the descending heights have root prefactors \(s^{L_i-j}\) \cite[Proposition~2.6]{Ko}.

\subsection{The two residues}

For \(j\ge i\) place \(\xi(w)\) at slice \(i\) and \(\eta(z)\) at slice \(j\), in that order also when the slices coincide. Commuting translations past multiplications gives
\[
 \langle N_i P_j\rangle=\operatorname{CT}_{w,z}
     N_i(w)P_j(z)K(z/w),\qquad
 K(h)=\frac{(1-h)(1-h/(ps))}{(1-h/p)(1-h/s)} ,
\]
where
\begin{equation}
\begin{split}
 P_j(z)&=\prod_{\beta<j,\,-}\frac{1-z/(sZ_\beta)}{1-z/Z_\beta}
          \prod_{\alpha\ge j,\,+}\frac{1-z/Z_\alpha}{1-z/(sZ_\alpha)},\\
 N_i(w)&=s^i\prod_{\beta<i,\,-}\frac{1-Z_\beta/(sw)}{1-Z_\beta/w}
          \prod_{\alpha\ge i,\,+}\frac{1-Z_\alpha/(cw)}{1-Z_\alpha/(pw)} .
\end{split}
 \label{eq:wick}
\end{equation}
The first expression absorbs \(s^{L_j}\) using the number of future pluses, and the second absorbs \(s^{-L_i}\) using the \(i+L_i\) past minuses. These formulas use the ordered two-current product, so shifts in a coefficient of a difference operator at another slice never need to be interchanged.

To evaluate the constant term, take \(u\) small with \(p,c\) generic and fixed. At a slice use circles between the scale of the future step parameters and that of the past steps, with the \(z\)-circle sufficiently smaller than the \(w\)-circle when the two slices coincide. All geometric expansions from the commutations converge on such circles, including both denominator expansions of \(K(z/w)\). Shrink the \(z\)-circle, then enlarge the \(w\)-circle. The zero pole contributes \(P_j(0)=1\); all other inner residues are
\[
 A_{j,\alpha}=\operatorname{Res}_{z=sZ_\alpha} P_j(z)\,dz/z
 \quad(\alpha\ge j,\,+).
\]
The outside constant is \(N_i(\infty)=s^i\); the other outer contributions are
\[
 B_{i,\beta}=-\operatorname{Res}_{w=Z_\beta} N_i(w)\,dw/w
 \quad(\beta<i,\,-).
\]
The poles of \(K\) are outside the shrinking contour in \(z\) and inside the expanding contour in \(w\). A pair of residues therefore carries just \(K(sZ_\alpha/Z_\beta)\). The calculation holds as a rational identity and hence coefficientwise in the specified expansion.

In \eqref{eq:diff} the constant inner term has zero forward difference and the constant outer term has zero backward difference. Extend the two residue arrays by zero when their step is absent. Each pair \(\beta<\alpha\) contributes, before its common factors, the finite sum
\begin{equation}
\sum_{i=\beta+1}^{\alpha}(A_i-A_{i+1})(B_i-sB_{i-1}),
 \qquad A_i=A_{i,\alpha},\quad B_i=B_{i,\beta}.
 \label{eq:sum}
\end{equation}
Put \(t_0=Z_\alpha/Z_\beta\) and \(a=Z_i/Z_\beta\). From \eqref{eq:wick},
\begin{equation}
(a-st_0)A_{i+1}=(a-t_0)A_i,\qquad
 (a-u)B_i=(a-su)B_{i-1}.
 \label{eq:rec}
\end{equation}
For the first relation either type of step gives \(P_{i+1}(z)/P_i(z)=(1-z/(sZ_i))/(1-z/Z_i)\). For the second, before evaluation at \(w=Z_\beta\), the ratio \(N_i/N_{i-1}\) is \((1-sw/Z_{i-1})/(1-w/Z_{i-1})\) at a preceding minus, and \((1-pw/Z_{i-1})/(1-cw/Z_{i-1})\) at a preceding plus. Inserting \(Z_i/Z_{i-1}=u\) or \(1/v\) gives \eqref{eq:rec} in both cases. The cleared relations hold at both endpoints, where \(A_{\alpha+1}=B_\beta=0\).

For \(t_0\ne u\), \eqref{eq:rec} implies
\begin{equation}
(A_i-A_{i+1})(B_i-sB_{i-1})
 =\frac{(s-1)t_0}{t_0-u}
       (A_{i+1}B_i-A_iB_{i-1}).
 \label{eq:tel}
\end{equation}
One can multiply out and use the cleared relations, including at \(a=u,t_0\). The sum telescopes to zero. Since \(t_0=u^{\alpha-\beta}c^{L_\alpha-L_\beta}\), for independent \(u,c\) equality \(t_0=u\) is possible exactly when \(\alpha=\beta+1\) and the intervening step is horizontal. Thus only the turns belonging to removable boxes survive.

At such a turn write \(Z_\alpha=c d_r\), \(Z_\beta=Z_\alpha/u\). Equation \eqref{eq:sum} consists of \(A_{\alpha,\alpha}B_{\alpha,\beta}\). The residues of the two adjacent steps in \eqref{eq:wick} are \(s-1\) and \(-s^\alpha(s-1)/s\), cancelling the multiplier in \eqref{eq:diff} up to \(s^\alpha\). The factors from the two other occurrences of the adjacent steps are \(R(u)^{-1}R(u/p)\); \(K(su)=R(u)/R(u/p)\) cancels these. For each earlier minus, the two factors give \(s^{-1}R(Z_\beta/Z_\gamma)/R(Z_\alpha/Z_\gamma)\); for each later plus they give \(sR(Z_\alpha/Z_\gamma)/R(cZ_\beta/Z_\gamma)\). There are \(\alpha+L_\alpha-1\) earlier minuses and \(L_\alpha-1\) later pluses, cancelling the remaining power \(s^\alpha\).

Along a horizontal run the factor from an earlier minus is
\[
 \frac{R(Z_\alpha/Z_{\gamma+1})}{R(Z_\alpha/Z_\gamma)} .
\]
Along a downward run the later plus factors give the inverse ratio. The two runs incident to \(r\) supply \(R(u)\) and \(R(v)\) after excluding the two adjacent steps. All other peaks supply \(R(d_r/d_h)\). Each valley, including the two outer endpoints, gives \(1/R(c d_r/d_a)\). This proves exactly \eqref{eq:corners} for the original descendant and finishes \eqref{eq:local}.

\section{Descent to the Picard family}\label{sec:descent}
\begin{proposition}\label{prop:descent}
Fix a diagram $\lambda$ and nonnegative integers $n,k$. If the two
affine vertices agree on all $x^a y^d$ with $0\le a\le n$ and
$0\le d\le k$, their contributions to the coefficient of
$q^{-b\sum h+n}t^k$ agree in the localized family group.
\end{proposition}
\subsection{The two effective divisors}
Write $U_A(H)$ for the push of $H|_A$ down a finite flat divisor,
$N_A(H)=\det U_A(H)/\det U_A(\OO)$, and $J_A=\det U_A(\OO)$.


Remove the common scalar
\[
 q^{-b\sum h}(-1)^{|\lambda|b}c^{b\sum(l+1/2)}.
\]
On the original side the empty remainder is $\prod E_c(Y^h s^{-l})$, where $Y=tL$ and $E_c(H)=\lambda_{-1}R\pi_*(cH)/\lambda_{-1}R\pi_*H$. On the rotated side it is $\prod E_s(q^h c^{-l})$. These are the empty Hilbert tangent pairs with the orientation of Lemma~\ref{lem:root}. In the rotated expression $R\pi_*(s q^h c^{-l})$ has underlying rank $b$ and the denominator rank $-b$, giving the empty product \eqref{eq:empty}.

For fixed excess $n$ on the original side take the double nested schemes $Z_{\lambda,d}$ of effective divisors $D_i$ ordered along all the arrows, with $\sum |D_i|=n$. They map finitely to $C^{(n)}$ by $D=\sum D_i$: the map is a restriction of the finite addition map from the product of symmetric powers. We use the scheme-theoretic incidence loci. The fully fixed virtual structure sheaf is their ordinary structure sheaf by \cite[Theorem~6.1]{MR}. This applies in the Picard family by pullback of the fully fixed theory: the box sheaves are the fixed divisor lines tensored by the indicated powers of $L$ and $K_C/L$. In a weight-zero morphism between two boxes the twists cancel against the identical coordinate monomial. They cancel in the section, arrow and commutativity terms and their differentials, so the relative fixed obstruction theory is pulled back from the reference curve calculation. On the rotated product the fixed schemes and their fixed theory are independent of $L$.

After removal of the empty scalar, the localized integrand can be encoded on $C^{(n)}\times C^{(k)}\times B$ before any family integration. The first color $n$ records the divisor $D$ above. The following description specifies the operations that define the second color; it also proves locality, which is needed for descent.

\begin{enumerate}\item Subtract the empty tangent before forming the relative determinant or fixed-weight Euler factor. The difference of each curve input is a sum of $V(A-B)-V$, where $V$ is a monomial line and $A,B$ are box divisors or zero. Its push is $U_A(V(A-B))-U_B(V)$ by the divisor sequences. The ordinary fixed structure sheaf is kept. In the inverse moving Euler use the full two-torus moving part. Fixed characters of the perfect obstruction theory are omitted there, also when several boxes have the same content.
\item The relative original determinant is the inverse determinant of the change of the Koszul polarization, with the prescribed scalar character and counting sign. The change has rank $n$. On the rotated side the modified root relative to the empty box module is $c^k(\det R\pi_*\Delta E)^{-1}N_{\sum D_i}(L/K_C)$ including the box characters in $E$, with the counting sign $(-1)^k$. This follows from the rank-$m$ determinant identity in the proof of Lemma~\ref{lem:root}; its root formula before splitting boxes applies to any rank-$m$ bundle. All relative determinants thus involve finite divisor data. Reflection of a negative counting weight brings its usual signed determinant. The $q$-degree of the rotated leading monomial stays $-b\sum h$ by the rank-derivative calculation in the proof of Lemma~\ref{lem:ring}, which also allows repeated contents.
\item After reflection, the remaining nonzero counting weights are positive. Restore the empty Euler factors from the beginning of this section at this step. Each positive-weight factor is the symmetric algebra of a signed sum of $t^lR\pi_*H$ with $l>0$ and $H$ a line of the form $V(A-B)$ or $V$. Duals are converted using relative curve duality. For each signed summand choose a marked divisor $M_\alpha$ of length $e_\alpha$ and weight $l_\alpha$, and send the tuple to $M=\sum l_\alpha M_\alpha$ of length $k$. A positive summand supplies $N_{M_\alpha}(H)$, a negative summand $N_{M_\alpha}(H)\det U_{M_\alpha}(1)[e_\alpha]$ in graded notation. Relative K\"unneth and the alternating covariants on the ordered cover give this formula with any base, since the alternating polynomials in each local coordinate are the Vandermonde times the symmetric polynomials. Push the product of these insertions and the preceding fixed factors along the two finite maps $(D_i),(M_\alpha)\mapsto(D,M)$. On the rotated side the stable-pair modifications give the second color $k$ and the positive $q$-weights the first color $n$. Coefficientwise pushforward from the two symmetric powers to $B$ recovers the original two family calculations.

\end{enumerate}

There may be factors of nonzero volume weight and zero counting weight. Before descent expand each such inverse Euler as a series in the volume double-cover character at zero, reflecting a negative weight first. With $n,k$ fixed there are finitely many fixed and marked length tuples. Each Laurent coefficient is a finite sum of operations on finite divisor bundles. This auxiliary expansion never requires division by a weight-zero Euler. It also avoids a nilpotence assumption on representation rings of stacks with stabilizers. After pullback to the projective family the expansion agrees with the usual rational axial localization, by the same line Euler identities. A Laurent polynomial with coefficients in $K^0(B)$ is detected by its expansion; localization at $1-\sigma^j$ is injective on these polynomials, including in the presence of additive torsion.

\subsection{Detection on colored punctual strata}

Use the atlas of two monic polynomials $f,g$ of degrees $n,k$. The arrows are changes of common \'etale coordinate modulo $(fg)^2$ separating the distinct support points, together with changes of frames for input lines. The target polynomials are the characteristic polynomials of the new coordinate on each quotient. Indeed substitution gives an isomorphism on the free rank-$2(n+k)$ algebras modulo $(fg)^2$: it does so on geometric fibers by separation and the derivative condition, then on families by Nakayama. Composition is substitution with the induced frame change. These arrows form a smooth groupoid (the allowed substitutions and frames are open conditions in vector bundles); denote its quotient stack by $\mathfrak D_{n,k}$. A family of pairs of effective divisors on a smooth curve maps to this stack. Locally on $C^{(n)}\times C^{(k)}\times B$ choose a common coordinate separating the support with nonvanishing derivative, and frames on the second neighborhood. The induced chart on the pair of symmetric powers is \'etale, as seen on completed local rings by coprime decomposition. With the smooth $B$ factor this gives flat pullback. Frame choices form a smooth torsor. The canonical bundle uses the differential of the coordinate on the divisor, which is determined modulo $fg$ by the substitution modulo $(fg)^2$. Constant nongeometric characters are recorded equivariantly. The nested conditions descend: divisibility of two monic factors means containment of their marked ideals, a closed condition compatible with base change and coordinate substitution. Thus their structure sheaves, also at a box with two predecessors, descend. Each relative bundle in the three steps uses a subdivisor of $D+M$, its conormal line, or the restriction of another divisor ideal; the product of any two such ideals contains the ideal of $2(D+M)$. For example $\mathcal O(-B)|_A$ is represented by $(f_B)/(f_Af_B)$. Norms $N_{M_\alpha}(V(A-B))$ are consequently defined on these data even at cross-color intersections.

In a stratum a support point now has type $(a_i,b_i)\ne(0,0)$. Its stabilizer consists of the disk jets and frame changes, with permutations of \emph{equal pairs}. The conormal at a single point has degrees $2,\ldots,a_i$ and $2,\ldots,b_i$, together with degree one when both colors occur. The latter is the difference of the two centers of the monic factors. The other coefficients come from centering each nonempty factor. Thus every central dilation weight of the conormal is positive. More precisely, the stabilizer at a configuration is, modulo its connected unipotent radical, the product of the dilation and frame tori at its points extended by permutations of identical types. Its integral coherent $K$-group is the representation group of this quotient, since the radical acts trivially on simple representations. Filter the stack by opens of at least $e$ support points. The new strata introduced at each stage are disjoint regular closed substacks there: the centered coefficients and, for two colors, the difference of centers form local regular parameters. By coherent localization a class vanishing on the preceding open is $i_*\alpha$ from the new closed strata; for these noetherian stacks coherent extension follows from \cite{Ry}. Its derived restriction is $\lambda_{-1}(N^\vee)\alpha$. The Euler has lowest simultaneous dilation term one with all others of positive degree, and so acts injectively even on the integral representation group with permutations. Induction detects the class by its derived restrictions. This argument is made before substituting $K_C$ for $p$.

On disjoint disks the nested ideals factor by coprime decomposition and so do the marked factors, with direct sums over choices of length allocations. The relative determinant of a rank-$d$ change occurs with shift $[d]$ from the counting sign, and every reflected Euler determinant likewise carries its rank shift. These are graded determinants under splitting. In particular the permutation actions are the tensor actions with Koszul signs; divisor norms split without an extra sign. Frame-independent scalar powers of the characters split because their exponents are additive for direct sums. The classes on the stratum are therefore products of the punctual classes with their full permutation actions. This description applies in each bounded-below axial expansion and includes the incidence of marks and fixed divisors.

For type $(a,b)$ its derived restriction at the single point is the corresponding affine coefficient of $A$ multiplied by $(p;p)_a(p;p)_b$, the cotangent Euler of the two symmetric powers. Indeed on that atlas the fixed stable-pair data are precisely the reverse plane partitions, and proper finite pushforward followed by affine localization gives the marked and nested expression above divided by these two Euler factors. The fiber character of the Poincar\'e line enters through $y=tL$ on the original side and the norm of $L$ on the rotated side, giving the same power $L^b$ at this point (here $b$ denotes the second local length). The full single-support comparison can therefore be done over the rational function field in independent $p,c$. After equality there, the integral torus representations in each axial coefficient agree. On a permutation cycle the supertrace is the punctual supercharacter with the torus entries multiplied along the cycle, so these equalities also give equality in the integral representation group including its permutation part. Applying stratum detection, flat pullback by the projective family's decorated divisors and proper pushforward proves the reduction.

\subsection{The marked curve and its index}
For $k\le1$ each support point has type $(a,0)$ or $(a,1)$.
Lemma~\ref{lem:boundary} and \eqref{eq:local} give every coefficient
required by Proposition~\ref{prop:descent}. Exchanging $(x,c)$
and $(y,s)$ gives the other strip, proving Corollary~\ref{cor:strip}.

For the explicit coefficient in Theorem~\ref{thm:mark}, use the
rotated fixed locus. A single modification has length one at a
removable box $r$. The moduli scheme of that mark is $C$, since
in families a degree-one effective divisor is a section and all the
framed monomial inclusions are forced. Put $W=d_r^{-1}$ and
$E_0=\sum_h d_h^{-1}$. The quotient increases by
$W(\OO(D)-1)$ at the universal mark $D$. The two divisor exact
sequences in Lemma~\ref{lem:root} give the tangent increment
\[
 K^{-1}H-cH^\vee,\qquad H=W[1-(1-u)(1-v)(E_0-W)^*],
\]
and the signed relative determinant $-cW^{-1}L$. Indeed the
determinant of cohomology contributes $(WK^{-1})^{-1}$ and the
pairing contributes $L/K$. On the disk $\det(\Delta H)=p^{-1}W$
so the signed factor $-s^{-1}/\det(\Delta H)$ agrees after removing
$L$. The mark has tangent $K^{-1}$, omitted from its projective
moving Euler. Hence multiplying the punctual contribution by
$1-p$ restores the tangent Euler \emph{before} setting $p=K$.
After this cancellation the insertion on the mark is the
$r$-summand of $(1-c)LG_\lambda$. Any remaining axial denominator
has nonzero $q$ degree: removable and addable contents strictly
interlace. Reflection to positive degree has a unit leading line.
Pushing to $B$ and multiplying by the empty factor \eqref{eq:empty}
proves \eqref{eq:mark} coefficientwise over the integral ring.

For the genus-two double the row and column give
$R(q)/R(q^2)$ and $R(c/q)/R((c/q)^2)$, respectively. Their empty
factors differ. At $q^{-2}t$ their sum is
\[
 D\bigl(L(c-1)(1+K+c^2/K+c^2)\bigr).
\]
Since $(K-1)^2=0$ on the curve, the first two coefficients in
$\delta$ are $4D(L)$ and $8D(L)-2D(L(K-1))$.
Thus the one-mark formula also detects a nonlinear volume term.

\begin{thebibliography}{Mon26}
\bibitem[Di12]{Di} D.-E.~Diaconescu, \emph{Moduli of ADHM sheaves and local Donaldson--Thomas theory}, J.~Geom. Phys. \textbf{62} (2012), 763--799. \href{https://doi.org/10.1016/j.geomphys.2011.12.018}{doi:10.1016/j.geomphys.2011.12.018}.
\bibitem[DS22]{DS} H.~Dinkins and A.~Smirnov, \emph{Quasimaps to zero-dimensional $A_\infty$-quiver varieties}, Int. Math. Res. Not. IMRN (2022), 1123--1153. \href{https://doi.org/10.1093/imrn/rnaa129}{doi:10.1093/imrn/rnaa129}. Equation numbers follow \href{https://arxiv.org/abs/1912.04834v2}{arXiv:1912.04834v2}.
\bibitem[FG10]{FG} B.~Fantechi and L.~G\"ottsche, \emph{Riemann--Roch theorems and elliptic genus for virtually smooth schemes}, Geom. Topol. \textbf{14} (2010), 83--115. \href{https://doi.org/10.2140/gt.2010.14.83}{doi:10.2140/gt.2010.14.83}.
\bibitem[Ko21]{Ko} S.~Koshida, \emph{Free field theory and observables of periodic Macdonald processes}, J.~Combin. Theory Ser.~A \textbf{182} (2021), 105473. \href{https://doi.org/10.1016/j.jcta.2021.105473}{doi:10.1016/j.jcta.2021.105473}.
\bibitem[Mon22]{M} S.~Monavari, \emph{Double nested Hilbert schemes and the local stable pairs theory of curves}, Compos. Math. \textbf{158} (2022), 1799--1849. \href{https://doi.org/10.1112/S0010437X22007606}{doi:10.1112/S0010437X22007606}.
\bibitem[Mon26]{MR} S.~Monavari, \emph{The refined local Donaldson--Thomas theory of curves}, Geom. Topol. \textbf{30} (2026), 3423--3484. \href{https://doi.org/10.2140/gt.2026.30.3423}{doi:10.2140/gt.2026.30.3423}.
\bibitem[NO16]{NO} N.~Nekrasov and A.~Okounkov, \emph{Membranes and sheaves}, Algebr. Geom. \textbf{3} (2016), 320--369. \href{https://doi.org/10.14231/AG-2016-015}{doi:10.14231/AG-2016-015}.
\bibitem[Qu18]{Q} F.~Qu, \emph{Virtual pullbacks in $K$-theory}, Ann. Inst. Fourier \textbf{68} (2018), 1609--1641. \href{https://doi.org/10.5802/aif.3194}{doi:10.5802/aif.3194}.
\bibitem[Ry16]{Ry} D.~Rydh, \emph{Approximation of sheaves on algebraic stacks}, Int. Math. Res. Not. IMRN (2016), 717--737. \href{https://doi.org/10.1093/imrn/rnv142}{doi:10.1093/imrn/rnv142}.
\bibitem[Sol63]{So} L.~Solomon, \emph{Invariants of finite reflection groups}, Nagoya Math.~J. \textbf{22} (1963), 57--64. \href{https://doi.org/10.1017/S0027763000011028}{doi:10.1017/S0027763000011028}.
\bibitem[Th22]{Th} R.~P.~Thomas, \emph{A $K$-theoretic Fulton class}, in \emph{Facets of Algebraic Geometry}, vol.~2, London Math. Soc. Lecture Note Ser. \textbf{473}, Cambridge Univ. Press, 2022, 367--380. \href{https://doi.org/10.1017/9781108877855.012}{doi:10.1017/9781108877855.012}.
\bibitem[W26]{W} B.~J.~Wuebben, \emph{The topological string on balanced local curves from two circle actions on a Calabi--Yau fivefold}, version of November 26, 2026. \href{https://doi.org/10.5281/zenodo.24549724}{doi:10.5281/zenodo.24549724}.
\end{thebibliography}

\end{document}
